% Abstract
\begin{abstract}
\noindent\textbf{Background.} The FlyWire FAFB v783 connectome provides the
complete wiring diagram of the adult female \emph{Drosophila} brain
(139,255 neurons; $\sim$3.7M thresholded connections at the pair level),
enabling neuron-level interrogation of directed network structure. Prior
work showed that high-degree neurons are disproportionately GABAergic,
but degree alone does not establish structural network control.

\noindent\textbf{Objective.} Test whether GABAergic identity predicts
neuron-level structural control impact --- the sensitivity of global
directed network efficiency to the removal of a single neuron --- beyond
what connectivity degree and degree-preserving network structure explain;
and characterize where the strongest structural chokepoints reside.

\noindent\textbf{Methods.} We defined the Control Impact Score,
$\mathrm{CIS}(i)=1-S(G-i)/S(G)$, with $S$ the sum of inverse directed
shortest-path lengths (freeze-$N$ convention), pre-registered 3,518
perturbation targets, computed CIS with a fixed-source-panel BFS estimator
(validated to $10^{-9}$ against exact computation), performed
pre-registered degree matching (GABA vs.\ ACh, total degree $\pm10\%$;
848 pairs), and compared the observed effect against 100 degree-preserving
null networks with exactly verified in/out degree sequences.

\noindent\textbf{Results.} Structural control impact is extremely skewed:
the median tested neuron accounts for $\sim$$10^{-5}$ of whole-brain
efficiency; the top neuron accounts for 2.40\%. The raw matched-pair
GABA--ACh effect was small (Cliff's $\delta=0.111$, Wilcoxon $p=0.0011$),
but a degree-controlled regression found no residual GABA association
($\beta_1=-0.025$, permutation $p=0.95$), and the observed statistic at
matched panel size ($\delta=0.098$) fell inside the degree-preserving null
ensemble (null $\delta=0.071\pm0.022$; empirical $p_\delta=0.109$;
median-difference $p=0.782$; $z=1.21$). In contrast, the top-50
chokepoints were strongly enriched for visual-centrifugal neurons
(21/50; 11.7$\times$ raw; 2.63$\times$ after degree matching, $z=5.3$,
$p=10^{-4}$; consistent at $K=25$ and $K=100$), and 13/50 individually
exceeded their degree-matched peers (10/13 visual-system), including a
low-degree octopaminergic centrifugal neuron at ${\sim}156\times$ its
peer-median impact.

\noindent\textbf{Conclusions.} The pre-registered GABA hypothesis is not
supported: the analysis does not provide evidence that GABAergic identity
independently predicts structural control impact after accounting for
degree and degree-preserving structure. Instead, removal-based
network-control sensitivity concentrates in visual-centrifugal
architecture, yielding a concrete, testable structural chokepoint map of
the adult fly brain. All claims concern structural network sensitivity,
not functional causality.
\end{abstract}
